\chapter{Blood in Urine}
\index{Blood in Urine}

\section*{Definition and Overview}
Blood in urine (haematuria) means that red blood cells are present in the urine. It may be visible to the naked eye (macroscopic or gross haematuria), turning the urine pink, red, or brown, or it may be invisible and detected only on a urine test (microscopic haematuria). The blood can originate at any point along the urinary tract -- the kidneys, the ureters that drain them, the bladder, the prostate (in men), or the urethra. Because a serious cause is possible, including infection, kidney stones, or cancer of the kidney or bladder, any blood in the urine warrants medical assessment, even when it appears to clear on its own.

\section*{Epidemiology}
Blood in urine is common. In younger women, urinary tract infections are by far the most frequent cause, while kidney stones also commonly cause bleeding and are more prevalent in men and in middle-aged people. The risk of bladder and kidney cancer rises with age, especially after 50 and in people who smoke, and visible blood in the urine of an older person always needs urgent investigation. Microscopic haematuria is found incidentally in a significant proportion of routine urine tests, and persistent microscopic haematuria in older adults carries a similar need for evaluation.

\section*{Aetiology and Pathophysiology}
Blood enters the urine at any point along the urinary tract:

\begin{itemize}
    \item \textbf{Urinary tract infection:} inflammation of the bladder (cystitis) or kidney (pyelonephritis) causes bleeding together with burning and frequency
    \item \textbf{Kidney stones:} stones irritate and scratch the lining of the kidney, ureter, or bladder, often producing severe colicky pain
    \item \textbf{Enlarged prostate:} in men, prostate disease can cause blood in the urine
    \item \textbf{Glomerulonephritis:} inflammation of the kidney's filtering units allows blood and protein to leak into the urine
    \item \textbf{Kidney, bladder, or ureteric cancer:} more likely in older people, smokers, and those exposed to certain industrial chemicals
    \item \textbf{Trauma:} injury to the kidney or bladder, including from contact sport, a fall, or a motor vehicle accident
    \item \textbf{Medicines:} anticoagulants and some other drugs can provoke bleeding
    \item \textbf{Strenuous exercise:} vigorous endurance exercise can rarely cause transient bleeding
\end{itemize}

\section*{Clinical Features}
\begin{itemize}
    \item Pink, red, or brown discolouration of the urine; a brown tint suggests older or kidney-derived blood
    \item Clots, which may be worm-like in shape if they formed in the ureter
    \item Burning or pain on urination, with urgency or frequency (infection)
    \item Severe, cramping pain in the flank or lower abdomen radiating to the groin (stones)
    \item Fever, chills, and loin pain (kidney infection)
    \item No symptoms at all in some cases -- blood found only on a routine test
    \item Unexplained weight loss or fatigue in older people (possible cancer)
    \item Swelling of the face or ankles in glomerulonephritis
\end{itemize}

\section*{Differential Diagnosis}
\begin{itemize}
    \item Urinary tract infection and kidney infection (pyelonephritis)
    \item Kidney and ureteric stones
    \item Prostate problems in men, including benign enlargement and prostate cancer
    \item Glomerulonephritis and other intrinsic kidney disease
    \item Kidney, bladder, or ureteric cancer
    \item Trauma to the kidney or bladder
    \item Menstrual blood contaminating a urine sample
    \item Beetroot, certain foods, or medicines that colour urine red without blood
    \item Dark urine from muscle breakdown (rhabdomyolysis) mistaken for blood
\end{itemize}

\section*{When to Seek Medical Care}
See a doctor for any visible blood in the urine, even if it clears quickly, and for an abnormal urine test showing blood. Seek prompt review if blood is accompanied by pain, fever, urinary symptoms, occurs after age 40, or occurs while taking blood thinners.

\textbf{Contact your local emergency services for:}
\begin{itemize}
    \item Inability to pass urine
    \item Heavy bleeding with large clots blocking the flow of urine
    \item High fever with shaking chills and loin pain
    \item Sudden severe flank or abdominal pain
    \item Dizziness, faintness, or collapse
\end{itemize}

\section*{Diagnosis and Investigations}
\begin{itemize}
    \item \textbf{History and examination:} timing and pattern of the bleeding, associated pain, current medicines, smoking history, and family history of kidney disease
    \item \textbf{Urine tests:} dipstick and microscopy to confirm blood and look for protein, casts, or infection
    \item \textbf{Urine culture:} to identify a urinary tract infection
    \item \textbf{Blood tests:} kidney function, electrolytes, and infection markers where relevant
    \item \textbf{Imaging:} ultrasound, or a CT scan of the kidneys and bladder (CT urogram), to look for stones, tumours, and structural abnormalities
    \item \textbf{Cystoscopy:} camera inspection of the bladder and urethra, especially in older people or when blood is visible
    \item \textbf{Repeat testing:} microscopic haematuria is often rechecked to confirm it is persistent before further investigation
\end{itemize}

\section*{Management}
\begin{itemize}
    \item \textbf{Infection:} antibiotics and plenty of fluids
    \item \textbf{Stones:} hydration and pain relief; smaller stones usually pass spontaneously, while larger ones may need lithotripsy or procedures to break up or remove them
    \item \textbf{Prostate disease:} treatment of benign enlargement or, if cancer is found, specialist urological care
    \item \textbf{Kidney disease:} specialist management of glomerulonephritis and related conditions
    \item \textbf{Cancer:} surgical removal and, where needed, other cancer treatments
    \item \textbf{Medicines:} blood thinners are reviewed but never stopped without medical advice
    \item \textbf{Exercise-related bleeding:} usually settles with rest; a doctor will check for other causes
    \item \textbf{Follow-up:} repeat urine tests to confirm the bleeding has resolved
\end{itemize}

\section*{Complications}
\begin{itemize}
    \item Anaemia if bleeding is heavy or prolonged
    \item Blockage of urine flow by a blood clot, causing pain and acute urinary retention
    \item Kidney damage if an underlying disease, such as glomerulonephritis, goes untreated
    \item A kidney or bladder cancer being found late if blood in urine is dismissed
    \item Recurrence of stones or infection
\end{itemize}

\section*{Prognosis and Outcomes}
The outlook depends on the cause. Most urinary tract infections and stones resolve fully with treatment, and exercise-related bleeding settles on its own. Persistent microscopic haematuria with a benign cause can be monitored safely. Bladder and kidney cancers have much better outcomes when caught early, which is why visible blood in urine -- especially in older people and smokers -- must always be investigated rather than assumed to be harmless.

\section*{Prevention and Self-Care}
\begin{itemize}
    \item Drink plenty of fluid, especially water, to keep the urine dilute and help prevent stones
    \item Treat urinary infections promptly and completely
    \item Do not smoke, and avoid harmful chemical exposures linked to bladder cancer
    \item Wear protective equipment for contact sport and seek care after any injury to the back or abdomen
    \item Have visible blood in the urine checked, even if it stops on its own
    \item Follow up abnormal routine urine tests as advised by your doctor
    \item Keep blood-thinning medicines under medical review and report any bleeding to the prescriber
\end{itemize}

\section*{Sources and Further Reading}
The following sources provide reliable, up-to-date information on blood in urine and urinary tract health:

\begin{itemize}
    \item NHS. \textit{NHS guidance on blood in urine and urinary symptoms.} https://www.nhs.uk/conditions/
    \item Mayo Clinic. \textit{Information on haematuria and kidney conditions.} https://www.mayoclinic.org/
    \item MedlinePlus. \textit{Blood in urine patient information.} https://medlineplus.gov/
    \item MSD Manual. \textit{Overview of urinary tract disorders and haematuria.} https://www.msdmanuals.com/
    \item NICE. \textit{Suspected cancer recognition and referral guidelines.} https://www.nice.org.uk/
    \item Centers for Disease Control and Prevention. \textit{Kidney disease and urinary tract infection information.} https://www.cdc.gov/
\end{itemize}
